\section{Introduction}
\label{sec:intro}

Parameter-Efficient Fine-Tuning (PEFT) methods, particularly Low-Rank Adaptation (LoRA) \cite{hu2021lora}, have revolutionized the way deep neural networks are adapted to downstream tasks. Rather than maintaining full-parameter copies of large pre-trained backbones, practitioners can serve a single base model paired with task-specific expert adapters. However, when serving these models on a highly heterogeneous stream where sequential samples require different skills, standard static weight merging \cite{taskvectors, tiesmerging} collapses due to parameter interference. This necessitates dynamic model ensembling, where routing layers dynamically blend the activations of specialized adapters on-the-fly at each layer \cite{shazeer2017outrageously, fedus2022switch}.

Existing dynamic ensembling methods typically operate in one of two regimes: stateless or stateful. Stateless routers, such as SABLE \cite{sable} or SPS-ZCA \cite{spszca}, compute ensembling coefficients at each layer based solely on the local activation's similarity to a set of pre-computed task centroids. While computationally efficient, these stateless methods suffer from a severe structural flaw in deep networks: representational noise. Because local activations fluctuate due to noise and cascading non-linearities, stateless routing coefficients exhibit massive layer-to-layer oscillations, a phenomenon we term \emph{routing jitter}. This jitter causes the network to blend incompatible expert projections in successive layers, initiating a cascade of representational drift that severely degrades the final accuracy \cite{qmerge}.

To resolve routing jitter, the state-of-the-art stateful router ChemMerge \cite{chemmerge} introduced a continuous physical framework. It models ensembling weights as chemical concentrations inside a continuous reactor, governed by non-equilibrium biochemical kinetics, Arrhenius reaction rates, and continuous-time Ordinary Differential Equations (ODEs) integrated via numerical solvers. By tracking expert concentrations as state variables across network layers, ChemMerge acts as a physical low-pass filter, dampening high-frequency routing oscillations and recovering significant performance.

\begin{figure}[t]
\centering
\includegraphics[width=0.42\textwidth]{performance_comparison.png}
\caption{Joint Mean Accuracy vs. Layer-to-Layer Routing Jitter (MSE) across ensembling methods. Our proposed \textbf{Momentum-Merge} outperforms the complex SOTA ChemMerge baseline while completely eliminating its biochemical ODE system, showing that a simple momentum update is highly stable and performant.}
\label{fig:performance_comparison}
\vspace{-0.5em}
\end{figure}

In this paper, we approach this continuous stateful formulation through the lens of Occam's razor. We ask a fundamental question: \emph{Are the biochemical kinetics, catalytic metaphors, and numerical ODE solvers truly necessary, or have we simply obscured a standard, classical smoothing filter under an overly convoluted physical apparatus?} 

By mathematically deconstructing ChemMerge, we show that under standard discretization, its continuous rate equations are equivalent to a state-dependent Exponential Moving Average (EMA). Following our research philosophy that \emph{simpler is strictly better if it matches complex performance}, we prune this unnecessary biochemical machinery. We propose \textbf{Momentum-Merge}, a training-free dynamic ensembling framework that stabilizes routing trajectories using a single, standard, constant EMA update across depth. Momentum-Merge requires exactly \emph{one} hyperparameter (the momentum coefficient $\beta$) and \emph{one} line of standard mathematical code, completely stripping away the chemical metaphors, Arrhenius reaction rates, and step-size discretization limits.

Despite its extreme simplicity, our empirical evaluation on a heterogeneous serving stream inside the Analytical Coordinate Sandbox (ICS) exposes a fundamental Accuracy-Stability trade-off: stateless calibrated routing (SABLE + Layer Centroids) achieves the highest joint accuracy (\textbf{77.24\%}), but suffers from massive layer-to-layer routing weight oscillations (routing jitter). Adding stateful smoothing acts as a low-pass filter that trades a small fraction of accuracy to achieve high routing stability. Under this trade-off, our basic Momentum-Merge achieves a joint mean accuracy of \textbf{74.85\%} (outperforming SABLE's best baseline accuracy of \textbf{74.76\%}) while reducing layer-to-layer routing jitter to \textbf{0.012860} (a \textbf{5.7$\times$} reduction over tuned SABLE). More significantly, our advanced Momentum-Merge variant reaches a joint mean accuracy of \textbf{74.98\%} (outperforming optimal SOTA ChemMerge of \textbf{74.71\%} by \textbf{+0.27\%} absolute) while recovering \textbf{92.41\%} of the performance of an Oracle expert execution and virtually eliminating routing trajectory oscillations, dropping routing jitter to an astonishing \textbf{0.000374}---representing a \textbf{195.7$\times$} reduction over tuned SABLE and a \textbf{41.1$\times$} reduction over tuned ChemMerge SOTA. We show how our single parameter $\beta$ controls the routing dynamics from stateless routing ($\beta=0$) to static uniform merging ($\beta=1$), peaking at $\beta = 0.60$. Our results prove that complex physical metaphors are redundant in dynamic model ensembling, paving the way for simpler, highly interpretable deep learning serving pipelines.

Our contributions are summarized as follows:
\begin{itemize}
    \item \textbf{Mathematical Deconstruction:} We prove that the chemical kinetics of ChemMerge are mathematically dual to a simple momentum filter under constant step discretization, exposing the physical metaphor as redundant complexity.
    \item \textbf{Minimalist Methodology:} We propose \textbf{Momentum-Merge}, which replaces the entire biochemical ODE solver with a single constant EMA equation, requiring only one hyperparameter ($\beta$) and zero system overhead.
    \item \textbf{Empirical Validation and Trade-off Mapping:} We evaluate Momentum-Merge in the Analytical Coordinate Sandbox, demonstrating that it outperforms SOTA while recovering 92.41\% of the Oracle ceiling. Symmetrically evaluating all methods reveals a fundamental Accuracy-Stability trade-off in dynamic model serving, mapping the boundaries of local expert plasticity and stateful trajectory smoothing.
    \item \textbf{Pareto Analysis:} We sweep the momentum parameter $\beta \in [0, 1]$, demonstrating the trade-off between local expert routing and representational smoothing, peaking at $\beta = 0.60$.
\end{itemize}
